\documentclass[10pt,a4paper]{report}
\usepackage[utf8]{inputenc}
\usepackage[T1]{fontenc}
\usepackage{natbib}
\usepackage[french, english]{babel}
\usepackage{lmodern, cmap}

\usepackage{graphicx, amssymb, amsmath, amsthm, mathtools, geometry, tikz}
\usepackage{multicol, hyperref, caption, subcaption, pgfplots}
\usepackage{mathrsfs}
\usepackage{paralist, microtype, url}
\usepackage[ruled,noend]{algorithm2e}
\usepackage{multirow}

\usepackage{tikz}
\usepackage{pgfplots}
\usetikzlibrary{3d, calc}
\usepackage{ulem} %pour le soulignement

\usepackage{lineno} %numérote lignes

\usepackage{booktabs}  % lignes plus fines et élégantes
\usepackage{array}
\usepackage{cleveref}


\makeatletter
% 1. Ne pas remettre à zéro les sections à chaque chapitre
\@addtoreset{section}{chapter}

% 2. Redéfinir les compteurs
\renewcommand{\thesection}{\arabic{section}}
\renewcommand{\thesubsection}{\thesection.\arabic{subsection}}
\renewcommand{\thesubsubsection}{\thesubsection.\arabic{subsubsection}}

% 3. (Optionnel) éviter que les chapitres apparaissent dans la numérotation ou la TOC
\renewcommand{\thechapter}{} % les chapitres n’ont plus de numéro
\makeatother



\SetKwComment{Comment}{/* }{ */}
\DeclareMathOperator{\ones}{ones}
\DeclareMathOperator{\diag}{diag}
\DeclareMathOperator{\som}{sum}
\DeclareMathOperator{\tr}{Tr}

%Nouvelles commandes et raccourcis
\renewcommand{\leq}{\leqslant}
\renewcommand{\geq}{\geqslant}
\newcommand{\ms}[1]{\mathscr{#1}}
\renewcommand{\eqref}[1]{eq. (\ref{#1})}
\newcommand{\figref}[1]{\textsc{Figure} \ref{#1}}
\newcommand{\tabref}[1]{\textsc{Table} \ref{#1}}


\newcommand{\Phx}{\Phi_x}
\newcommand{\Phy}{\Phi_y}
\newcommand{\X}{\mathcal{X}}
\newcommand{\Y}{\mathcal{Y}}
\newcommand{\C}{\mathcal{C}}
\newcommand{\I}{\mathcal{I}}
\newcommand{\U}{\mathcal{U}}
\newcommand{\R}{\mathbb{R}}
\newcommand{\N}{\mathbb{N}}
\newcommand{\E}{\mathbb{E}}
\renewcommand{\P}{\mathbb{P}}
\renewcommand{\H}{\mathscr{H}}
\newcommand{\G}{\mathscr{G}}
\newcommand{\fat}{f_\theta^a}
\newcommand{\fbt}{f_\theta^b}
\newcommand{\norm}[1]{\left\|#1\right\|}
\newcommand{\nhs}[1]{\left\|#1\right\|_\text{HS}}
\newcommand{\scal}[2]{\left\langle #1 \vert #2 \right\rangle}

\pgfplotsset{compat=1.17}


\usepackage{libertinus}







\begin{document}


\begin{titlepage}
    \centering

    \noindent\rule{\textwidth}{1mm}\\ 
    {\Large \bf Dimension reduction and manifold learning \par}
    \vspace{0.4cm}
    {\huge\sc Minimax Estimation of distances on a surface and minimax manifold learning in the isometric-to-convex setting\par}
    \noindent\rule{\textwidth}{1mm}\\ 
    \vspace{2cm}
    
    % Auteur 1
    {\sc{Abel Douzal} \par}
    {\textit{École normale supérieure PSL} \par}
    {\href{mailto:abel.douzal@ens.fr}{abel.douzal@ens.fr}}
    \vspace{2cm}
    
    % Auteur 2
    {\sc{Ruben Ifrah} \par}
    {\textit{École polytechnique} \par}
    {\href{mailto:ruben.ifrah@polytechnique.edu}{ruben.ifrah@polytechnique.edu}}
    \vspace{2cm}

    % Auteur 2
    {\sc{Pierre Cornilleau} \par}
    {\textit{Dauphine PSL} \par}
    {\href{mailto:pierre.cornilleau@dauphine.eu}{pierre.cornilleau@dauphine.eu}}
    \vspace{5cm}

    %\begin{center}
    %    \includegraphics[width=0.4\textwidth]{Images/dauphine.png}
    %    \hspace{2cm}
    %    \includegraphics[width=0.2\textwidth]{Images/Logo_École_normale_supérieure_-_PSL_(ENS-PSL).svg.png}
    %\end{center}


    \vfill
    % Date
    {\large \today\par}
\end{titlepage}

\newgeometry{top=3cm, bottom=3cm, left=1.5cm, right=1.5cm}
\setlength{\columnsep}{1cm}
\tableofcontents





%\linenumbers %numérote lignes



\section{Introduction}



\section{Conclusion}




\end{document}